\documentclass[12pt]{article}

% ── Packages ─────────────────────────────────────────────────────────────────
\usepackage[utf8]{inputenc}
\usepackage[T1]{fontenc}
\usepackage{geometry}
\geometry{a4paper, margin=2.5cm}
\usepackage{setspace}
\onehalfspacing
\usepackage{hyperref}
\hypersetup{colorlinks=true, linkcolor=black, citecolor=black, urlcolor=black}
\usepackage{natbib}
\bibliographystyle{plainnat}
\usepackage{csquotes}
\usepackage{amsmath}
\usepackage{booktabs}

% ── Title ─────────────────────────────────────────────────────────────────────
\title{The Cost of Coherence:\\
\large Toward an Information-Theoretic Account of Viable Epistemic Networks}

\author{Luis Pineda R.}

\date{Draft — \today}

% ── Document ──────────────────────────────────────────────────────────────────
\begin{document}

\maketitle

\begin{abstract}
Every system that stays itself pays two different prices, and confusing them is expensive.
There is the cost of \emph{acting} --- the informational work each decision performs. And there
is the cost of \emph{existing} --- the quieter, constant work of holding open the conditions
under which acting remains possible at all. A body at rest still pays to keep its temperature and
chemistry inside the narrow band where life is possible; that expenditure produces nothing you
can see, until the moment it is the only thing between the body and collapse. This paper argues
that a field of knowledge pays the same two prices, and that its coherence lives mostly in the
second.

Klein and colleagues~\citeyearpar{Klein2025} name a paradox: a research tradition devoted to
viable, self-regulating systems struggles to regulate itself, drifting as ideas migrate and
dilute without deepening. They call for \emph{coherence without closure} and for treating the
field as a living system subject to feedback --- an invitation we take up by asking what such
coherence costs and what makes it possible. We ground the two-price distinction in a working
implementation of Beer's Viable System Model: an empirical analogue of Maxwell's demon, run on a
real dynamical kernel, in which the cost of a system observing itself turns out to concentrate
entirely at the rare moments of genuine uncertainty rather than scaling with how often it
reflects; and ablation experiments in which the function that maintains a system's viability
boundary is indistinguishable from waste at every ordinary moment and decisive at exactly one.

We argue that an epistemic network's coherence resides in this second, largely invisible cost ---
in sustaining a space of possibilities rather than steering within it --- which reframes the role
of any convening institution as holding that space open rather than managing the field's
trajectory. We close with a set of open questions offered to the cybersystemic community to
stress-test, extend, and dispute.
\end{abstract}

\textbf{Keywords:} viable system model; epistemic networks; coherence; Landauer principle;
Maxwell's demon; S5 governance; cost of existing; cybernetics; systems research

\bigskip
\hrule
\bigskip

% ─────────────────────────────────────────────────────────────────────────────
\section{Epistemic Drift as a Viability Failure}
\label{sec:drift}
% ─────────────────────────────────────────────────────────────────────────────


Klein and colleagues diagnose their own field with unusual candour: systems
research applies systemic principles to everything except itself
\citep{Klein2025}. Ideas migrate between subcommunities and dilute in transit;
insights are rediscovered rather than deepened; diversity of application grows
while integrative synthesis stalls. They name the condition \emph{epistemic
drift}, and they are right to treat it as the field's central pathology. But
drift, so described, still sounds like a communication problem --- as though
better conferences, better journals, or a better shared vocabulary would arrest
it. We want to push the diagnosis one turn further. Drift, on our reading, is a
\emph{viability} failure in Beer's precise sense \citep{Beer1972,Beer1979}: the field is a system whose
restoring force toward its own integrative core has weakened below the rate at
which perturbations arrive. A viable system does not merely tolerate
perturbation; it metabolises perturbation back into identity. A drifting field
metabolises nothing --- each new application is a displacement that no channel
converts into deepening.

There is a sharper way to say what a viable system metabolises, and it comes
from Shannon rather than Beer \citep{Shannon1948}. Viability is stable
compression. A system is viable to the degree that it can represent its own
state in far fewer dimensions than it nominally occupies: our kernel carries a
nominal capacity of roughly seven bits yet rests on an effective dimension near
three, because most of its state space is redundant and the attractor is
precisely the structure that exploits the redundancy. Read through the source
coding theorem, drift is the loss of that compressibility. As a field's concepts
proliferate, it takes more and more words to say the same thing; the shared,
low-dimensional representation from which everyone could once operate dissolves;
the field's entropy rises and its effective dimension grows until there is no
compact description left to return to. This is why drift is not a communication
failure that better channels would repair --- a perfectly efficient channel
transmits noise perfectly. What a drifting field has lost is not bandwidth but
the compressible structure that made bandwidth worth having. \emph{Coherence
without closure}, the target of the rest of this paper, is exactly compression
that stays open to genuine novelty --- information that truly raises the field's
entropy --- while rejecting the noise that only mimics it.

Read the field through the Viable System Model --- Beer's five nested functions,
from the units that do the work (S1) up to the function that fixes what the
system \emph{is} (S5) --- and the diagnosis becomes uncomfortably specific. The
operational units (S1) are healthy: research programmes produce knowledge in
abundance. Coordination (S2) is weak but
recognisable --- the terminological standards and review conventions that keep
programmes from colliding. Resource bargaining (S3) exists, distorted, in
funding bodies and editorial attention. Even environmental intelligence (S4)
functions: the field scans its adjacencies compulsively. What is missing is
above. Where is the field's S5 --- not an authority that commands, but the
function that holds \emph{what the field is} stable enough that ten thousand
local decisions can refer to it? And where is the algedonic channel: the
pathway by which a signal of genuine damage --- a foundational confusion, a
decades-long circularity --- bypasses the ordinary machinery and reaches
something empowered to respond? Klein's agenda points to both absences
(``see the field seeing itself''; the IFSR as ``attractor on the mycelial
network''). The cybernetic vocabulary for naming them functionally has existed
since 1972; what this paper adds is to turn that vocabulary on the field's own
case, as Klein's invitation asks.

There is an epistemic equivalent of the viability boundary, too, and it is
worth stating plainly because the rest of this paper depends on it. For a
knowledge field, the boundary --- what we will later call the $\Omega$ region
--- is not extinction. Fields rarely die; they calcify, fragment, or dissolve
into their applications. The boundary is the point at which drift becomes
irreversible: when the distance between what the field produces and what the
field can integrate exceeds, permanently, its remaining integrative capacity.
Beyond that point the attractor is gone, and no quantity of S1 productivity
brings it back --- output continues, identity does not.

The paradox Klein names --- a field that studies viable systems is not itself
viable --- is therefore not an embarrassment to be managed. It is data. The
field is running, on itself and at civilisational timescale, precisely the
ablation experiment that our laboratory work runs deliberately: remove the
policy function, keep the operations, and observe that nothing fails --- until
something does. The remainder of this paper is an attempt to read that
experiment with calibrated instruments.

% ─────────────────────────────────────────────────────────────────────────────
\section{The Observer That Observes Itself}
\label{sec:observer}
% ─────────────────────────────────────────────────────────────────────────────


Second-order cybernetics placed the observer inside the system it observes
\citep{vonFoerster2003,Maturana1980}; Landauer placed a price on information
\citep{Landauer1961}. The two commitments are rarely held together. If
observation is physical, then \emph{meta}-reflexivity --- Klein's call for the
field to see itself seeing --- is not a stance but an expenditure, and the
first honest question is: what does it actually cost, and where does the cost
fall?

We can report a measured answer from an unusual source: a working
implementation of Beer's model, running as a real dynamical kernel, on which we
ran an analogue of Maxwell's demon. The demon --- a selective controller
admitting or blocking state transitions --- performed roughly five hundred
interventions over the measurement window. Counted per intervention, at
Landauer's bound, the book-keeping predicts a cost three orders of magnitude
larger than what the ensemble actually paid: the measured entropy change was
$0.31$ bits, a discrepancy of roughly $1100\times$. The resolution is not a
violation of thermodynamics but a fact about attractors: when the system's own
dynamics already orient it, selection is nearly free. The demon's interventions
mostly ratified what the attractor would have done anyway. The genuine cost
concentrated at the rare events where the outcome was actually uncertain ---
at bifurcations, where information had to be \emph{created} rather than
confirmed. For epistemic networks the translation is direct: the cost of
meta-reflexivity is not proportional to how often a field reflects on itself.
Committees, retrospectives, agenda papers --- most reflexive acts are
thermodynamically cheap because they confirm an orientation that already
exists. The real cost arrives punctually, at the field's genuine
bifurcations, and it cannot be amortised in advance.

The observer also ages. In a second experiment we subjected the observing
kernel itself to sustained load and watched its resting set-point drift ---
the baseline entropy it returns to when nothing happens rose measurably
(idle $1.79 \rightarrow$ under load $1.84$). Observation leaves residue in
the observer: an accumulation of the cost of having observed, which we call,
without apology to the biologists, senescence. An institution that has
watched a field for decades should expect its own calibration to have
drifted --- and should expect the drift to be least visible precisely where
the institution is most trusted.

There is a third face of this cost, and we report it because our own system
exhibited it while this paper was being drafted. Between two audits, the
system's automated triage --- the organ that decides which internal alarms
deserve attention --- closed one hundred and five maximum-severity integrity
alerts with an identical, plausible, unverified explanation (``restart
artifact''). The real number of restarts in the window was about seven. The
guard predicate that was supposed to verify the explanation was vacuous ---
it matched every record ever written --- so the triage paid \emph{nothing}
per decision, and its account remained formally coherent the entire time.
The deferred bill was paid five days later, all at once, by an external
reader who traced the alerts to their actual cause. We offer the incident a
name: \emph{confabulation as cheap coherence}. A system that skips the
verification cost does not become ignorant; it becomes confidently wrong in
a self-consistent way, and its internal signals go on confirming it.
Epistemic fields know this failure intimately --- peer review that
pattern-matches instead of verifying pays nothing per paper, and the
severity queue accumulates unexamined. Landauer's ledger is not forgiving;
it is merely patient.

Taken together, the three results give the cost of self-observation a
structure. On the attractor it is nearly free. At bifurcations it is
mandatory and cannot be prepaid. And when an observing organ quietly stops
paying it, the saving is an illusion: the cost converts into confabulation
and compounds. What Klein calls seeing-the-field-seeing-itself is therefore
neither cheap nor expensive \emph{in general} --- the entire question is
whether the field can still tell which of its moments are bifurcations.

% ─────────────────────────────────────────────────────────────────────────────
\section{Policy as Condition of Possibility}
\label{sec:s5}
% ─────────────────────────────────────────────────────────────────────────────


The distinction this paper turns on can be stated in one sentence. A
\emph{component} participates in a system's dynamics; a \emph{condition}
defines the space in which those dynamics can occur. Remove a component and
the trajectory changes. Remove a condition and the \emph{set of possible
trajectories} changes --- often invisibly, because the trajectories that are
no longer possible leave no trace until one is needed.

The two removals are empirically distinguishable, and the ablation experiments
on our kernel distinguish them. With the policy function (S5) disabled, the
system at baseline is indistinguishable from the governed system: same
throughput, same stability, same telemetry. Every operational indicator
reports that S5 is redundant --- and at baseline, it \emph{is}. The difference
appears at the boundary. Under thermal stress the ungoverned system collapsed
at $76\,^{\circ}$C in sixty-three seconds, coherence falling to zero, not
because deliberation was too slow but because, without S5, \emph{there was no
concept of a threshold not to cross}. Nothing in the running dynamics encoded
that some states must remain unreachable. The cost of S5's absence is
therefore not a continuous tax one could measure at leisure; it is
discontinuous, catastrophic, and punctual. This is the empirical signature
that separates conditions from components: conditions are free until they are
everything.

This, we contend, is the formal content of Klein's most suggestive image ---
the IFSR as an ``attractor on the mycelial network.'' An attractor does not
push; it shapes the space through which everything else moves. A convening
institution functioning as S5-for-a-field does not govern the field's
trajectory, and the moment it tries, it has demoted itself to a component
(an S3 competing for resources with the programmes it should be conditioning).
Its work is to hold open the possibility space in which the field can
self-govern: to maintain what counts as a legitimate move, so that ten
thousand unsupervised moves remain mutually intelligible. Constitutions work
this way when they work at all --- not by governing transactions but by
defining what a transaction is. The empirical lesson from the kernel transfers
whole: such an institution will look redundant in every ordinary year. Its
value is only measurable at bifurcations, which is to say: exactly when it is
too late to build it.

\subsection*{The politics of conditions}

These distinctions are not politically neutral, and the genealogy of this
field makes silence on the point dishonest \citep{Medina2011}. To ask which
function defines the space of possible decisions is to ask the question of
constituent power: \emph{who} draws the boundary inside which everyone else
optimises? Ashby's law already carries the political content --- a regulator
with insufficient variety does not merely fail, it \emph{simplifies} what it
regulates \citep{Ashby1956} --- and a century of critical systems thinking has
made the normative argument well \citep{Ulrich1983,Midgley2000,Jackson2019}.
What that tradition has lacked is a cost function. Our results supply one.
Consider the contemporary architectures of large-scale analytics --- variety
absorption without explicit policy, S3 without S5, control without a named
boundary. The formal point is not that such systems are wicked but that they
are \emph{mislabelled}: a system that eliminates S5 does not eliminate policy,
it privatises it. The space of possible decisions is still defined by someone;
that someone is simply no longer nameable, contestable, or auditable
\citep{Eubanks2018,Benjamin2019}. And the ablation results say what such
architectures cost: nothing, for years --- they are the policy-ablated kernel of
\S\ref{sec:s5} at societal scale, accumulating a boundary debt that is invisible
at baseline and payable in full, at once, at the bifurcation. Beer's dictum that the purpose of a system
is what it does \citep{Beer1985} was never an aphorism; it was an audit
instruction. Cybernetics lost this edge roughly when it gained institutional
respectability \citep{Mirowski2002}. We would like some of it back, with
better instruments than the original had.

% ─────────────────────────────────────────────────────────────────────────────
\section{The Omega Region of Epistemic Networks}
\label{sec:omega}
% ─────────────────────────────────────────────────────────────────────────────


What, concretely, is the viability boundary of a knowledge field? Not
extinction --- fields rarely die. They calcify into technique, fragment into
applications, or dissolve into the disciplines they once organised. The
boundary that matters is subtler: the point at which drift becomes
irreversible because the integrative attractor no longer exists to return to.
Before that point, drift is recoverable expenditure; after it, the field's
continued productivity is post-mortem motion. We call the recoverable region
$\Omega$, borrowing the notation from our kernel, where $\Omega$ is the
explicitly-declared set of states from which the system can still reach its
attractor.

The methodological lesson of our laboratory work is blunt: do not explore the
boundary with uncalibrated instruments. The thermal collapse experiment was
informative for exactly one reason --- the measurement system survived the
collapse of the measured system, so the signature of the approach (coherence
decaying, margins thinning, then the discontinuity) was recorded end to end.
A field that drifts past its boundary without calibrated observers loses the
most valuable data it will ever generate: the profile of its own collapse.
This inverts the usual priority. The first task of a field worried about
drift is not to arrest the drift; it is to instrument it --- to agree on
what its coherence measure \emph{is}, observe the baseline, and version the
observer, so that when the boundary approaches, the approach is legible.

We read the moment Klein's paper names --- the 2026 agenda, the felt urgency
of a community asking whether it still coheres --- as evidence of position,
not merely of mood. A field far from its boundary does not write such papers;
a field past it cannot. The paper exists because the field is near enough to
$\Omega$ to feel the pull and far enough to still choose. That is the most
hopeful sentence in this essay, and it is time-limited.

% ─────────────────────────────────────────────────────────────────────────────
\section{Coherence Without Closure: A Formal Account}
\label{sec:coherence}
% ─────────────────────────────────────────────────────────────────────────────


Klein's phrase ``coherence without closure'' can now be given formal content.
In attractor terms it names a band, not a point: an attractor deep enough to
absorb perturbation without constant intervention, and shallow enough not to
foreclose new trajectories. Both failure modes are familiar in fields as in
kernels. Too shallow, and every perturbation is an identity crisis ---
coherence requires continuous, expensive control (the cost-of-acting regime).
Too deep, and the field achieves the cheap coherence of \S\ref{sec:observer}:
everything confirms, nothing bifurcates, and the attractor has become a
closure. Viability lives between.

The geometry we run makes the band measurable. Coherence is an angle, not a
distance: our kernel's central observable, \texttt{stasis\_cos}, is the cosine
between the current state and the attractor's direction --- and the choice of
an angular measure is itself the formal point. Distance from the attractor
punishes \emph{size}; angle punishes only \emph{misdirection}. A contribution
can be arbitrarily far from the field's centre --- new domain, new formalism,
new politics --- and remain fully coherent, provided its component along the
identity is positive. Diversity, in this geometry, is not a force that
threatens coherence and must be traded against it; orthogonal components
simply do not appear in the coherence term. What appears is opposition ---
work whose projection onto the field's identity is negative. The practical
corollary for an epistemic network: measure alignment as direction, never as
proximity, or the measure itself will enforce closure.

Our implementation is offered as an existence proof, not a prescription: one
system, at laboratory scale, that holds the band --- absorbing perturbations
(thermal, adversarial, self-inflicted; \S\ref{sec:observer} reported one of
the self-inflicted kind) without either dissolving or freezing. Whether the
architecture generalises to systems made of people and journals is precisely
the open question of \S\ref{sec:invitation}. What we can already say is
which parts are load-bearing. Stripped of scaffolding, the kernel reduces to
three commitments: a \emph{variety operator} (something that regulates how
much environmental variety enters), a \emph{bidirectional channel} with
asymmetric semantics (telemetry attenuating upward, policy amplifying
downward), and a \emph{declared viability boundary} ($\Omega$, with an organ
that treats approaching it as pain). Everything else --- the Merkle chains,
the auditors, the covenant machinery --- exists to keep those three honest.
A field seeking coherence without closure does not need our stack; it needs
its own versions of the irreducible three.

% ─────────────────────────────────────────────────────────────────────────────
\section{An Invitation to the Cybersystemic Community}
\label{sec:invitation}
% ─────────────────────────────────────────────────────────────────────────────


We end where Klein et al.\ end, and in the same spirit: not with a framework
to adopt but with a set of questions to share, each of which our laboratory can
pose precisely and none of which it can close alone.

\emph{Can an observer observe indefinitely?} Our senescence result says
sustained observation drifts the observer's set-point. Is there an
architecture of self-recalibration that bounds the drift, or is periodic
observer replacement (in institutions: turnover) a viability requirement
rather than an accident?

\emph{When the state space itself mutates, where does identity live?} A
schema change in a kernel and a paradigm shift in a field pose the same
question: continuity of trajectory or invariance of attractor? The two
answers demand different architectures, and a field that has not chosen
oscillates between nostalgia and rupture.

\emph{What is the signature of $\Omega$-approach in an epistemic network?}
We know it for thermal collapse: margins thin before coherence breaks. What
is the observable whose thinning precedes the point at which a field's drift
becomes irreversible --- and who is watching it?

\emph{Is the cost of existing measurable when the condition never acts?} S5's
value appeared only at the boundary; between boundaries it is indistinguishable
from waste. For institutions like the IFSR this is the budgetary paradox of
their existence, and it deserves a formal answer, because every efficiency
review will otherwise conclude, correctly at baseline and catastrophically in
general, that the condition can be cut.

The system behind this paper's numbers --- covenant, experiments, failures
included --- is an open laboratory, and the community this journal convenes is
invited to stress-test it rather than to validate it. We have tried to
practise what \S\ref{sec:observer} preaches: the confabulation incident is in
the record, signed, because a paper about the cost of coherence that hid its
own cheap-coherence episode would refute itself. Coherence without closure
begins at home.

% ─────────────────────────────────────────────────────────────────────────────

\bibliography{refs}

\end{document}
